\documentclass[11pt]{article}
\usepackage[margin=1in]{geometry}
\usepackage{amsmath,amssymb}
\usepackage{xurl}
\usepackage[hidelinks]{hyperref}
% Shared commands for notes in docs/paper-gaps/.

\DeclareUnicodeCharacter{2080}{\ifmmode{}_{0}\else\textsubscript{0}\fi}
\DeclareUnicodeCharacter{2081}{\ifmmode{}_{1}\else\textsubscript{1}\fi}
\DeclareUnicodeCharacter{2082}{\ifmmode{}_{2}\else\textsubscript{2}\fi}
\DeclareUnicodeCharacter{2099}{\ifmmode{}_{n}\else\textsubscript{n}\fi}

\newcommand{\Lean}{\textsc{Lean}}
\ExplSyntaxOn
\NewDocumentCommand{\leanid}{m}{
  \ifmmode
    \textsf{\detokenize{#1}}
  \else
    \group_begin:
    \urlstyle{sf}
    \tl_set:Nn \l_tmpa_tl {#1}
    \tl_replace_all:Nnn \l_tmpa_tl {\_} {_}
    \exp_args:NV \path \l_tmpa_tl
    \group_end:
  \fi
}
\ExplSyntaxOff
\providecommand{\tr}{\operatorname{tr}}
\newcommand{\tnleanIssueBase}{https://github.com/LionSR/TNLean/issues/}
\newcommand{\tnleanPrBase}{https://github.com/LionSR/TNLean/pull/}
\newcommand{\tnleanDocBase}{https://sirui-lu.com/TNLean/docs/}
\newcommand{\ghissue}[1]{\href{\tnleanIssueBase#1}{GitHub issue \##1}}
\newcommand{\ghpr}[1]{\href{\tnleanPrBase#1}{PR \##1}}
\newcommand{\doclink}[1]{\href{\tnleanDocBase#1.html}{\path{#1}}}

% Dirac notation and the identity operator, as in blueprint/src/macros/common.tex.
% \providecommand keeps note-local definitions authoritative.
\usepackage{dsfont}
\providecommand{\ket}[1]{|#1\rangle}
\providecommand{\bra}[1]{\langle #1|}
\providecommand{\braket}[2]{\langle #1 | #2 \rangle}
\providecommand{\ketbra}[2]{|#1\rangle\!\langle #2|}
\providecommand{\Id}{\mathds{1}}
\providecommand{\one}{\mathds{1}}
\providecommand{\C}{\mathbb{C}}
\providecommand{\MN}[1]{M_{#1}(\C)}
\providecommand{\MD}{M_D(\C)}

% Verdict marker: \gapnote{<kind>}{<status>}. Kind is the policy
% classification (clarification, local-correction, scope-restriction,
% unfaithful, false-source, open-gap); status is open, wip (elimination
% underway), resolved, or historical. Machine-read by the site index;
% prints a verdict line.
\newcommand{\gapnote}[2]{\par\noindent\textbf{Verdict:} #1 (#2).\par}

\title{Nonempty parties in distributed density compression}
\author{}
\date{2026-10-08}
\begin{document}
\maketitle
\gapnote{local-correction}{resolved}

Theorem~5.2 of the September 24, 2026 polynomial-PEPS manuscript concerns
compression to an operator tensor network with one node per party
\cite[Theorem~5.2]{OpenAI2026PolynomialPEPS}.
The statement does not explicitly require the party set to be nonempty.
The formalization adopts this requirement once, in the definition of a
\emph{distributed construction}. It is the intended convention needed by
the scalar-absorption step of the proof, rather than an assertion that the
printed statement includes this qualifier.

\paragraph{The scalar boundary case.}
The preliminaries state that empty tensor products have dimension one
(\texttt{01-preliminaries.tex}, lines 5--7)
\cite[preliminaries]{OpenAI2026PolynomialPEPS}.
Accordingly, a circuit on no parties can still act on $\mathbb C$ by a scalar
contraction $c$. Starting from the empty product input $1$, its physical
density is the scalar $|c|^2$. An operator tensor network with no nodes or
links contracts to the empty product $1$. These are different when, for
example, $c=1/2$; their absolute difference is $3/4$.
Thus the literal zero-party reading does not imply approximation to
arbitrarily small error by a network having exactly one node per party.
No zero-party compression theorem is asserted.

\paragraph{The adopted convention.}
A distributed construction has a finite nonempty set of parties.
The original memories are finite-dimensional, the initial vector is a
product of normalized party vectors, and the circuit retains the private
output registers until the final trace, exactly as in
\texttt{04-compression.tex}, lines 17--30 and 137--151
\cite[Theorem~5.2]{OpenAI2026PolynomialPEPS}.
No additional bound is imposed on private dimensions.
The network construction chooses local tensors at existing parties and
absorbs scalar coefficients and averaging factors into them
(lines 565--578) \cite[proof of Theorem~5.2]{OpenAI2026PolynomialPEPS}.
In particular, the scalar contributed by a gate with no
participating parties can be absorbed at any existing party, without
introducing a node or a virtual link.

The general circuit definition remains unrestricted: gates with no
participating parties, their chronological occurrences, and their scalar
contributions are all retained. Only the definition of a distributed construction carries the
nonempty-party convention; it is not repeated as a hypothesis of each
subsequent theorem. The spatial families used for
square grids of side length $L\geq2$ already have a nonempty party set,
so this convention does not change those applications.

\paragraph{Source.}
The cited source is the manuscript at
\href{https://github.com/openai/math/tree/adc7f1241b42e322a6451854ab7e4b4c146bf78a}{\texttt{openai/math@adc7f1241b42}},
\emph{Polynomial PEPS approximation of gapped square-grid ground states},
the preliminaries and the section on distributed density compression
\cite{OpenAI2026PolynomialPEPS}.

\bibliographystyle{alpha}
\bibliography{references}

\end{document}
